The cepheides are variable stars, whose luminosities vary due to volume
oscillations. The period of the oscillations of a cepheide star is:
\[ P = 2\pi R \sqrt{\frac{R}{KM}}, \]
where $R$ is the radius of the cepheide; $M$ is the mass of the cepheid
(constant during oscillation); $R = R(t)$; $P = P(t)$.

Demonstrate that the absolute magnitude of the cepheide, $M_{\mathrm{cef}}$,
depends on the period of the cepheide's pulsation $P$ according to the
following relation:
\[ M_{\mathrm{cef}} = -2.5^{\mathrm{m}} \cdot \log k
   - \left(\frac{10}{3}\right)^{\mathrm{m}} \cdot \log P, \]
where $k$ is a constant; $P = P(t)$;
$M_{\mathrm{cef}} = M_{\mathrm{cef}}(t)$.
